\section{Scalar words and exact itineraries}\label{sec:itineraries}

Writing $T^iP=(x_i,x_{i+1},x_{i+2})$ turns the map into the
two-sided scalar recurrence
\begin{equation}\label{eq:recurrence}
 x_{i+3}=x_{i+1}x_{i+2}-x_i,
 \qquad x_{i-1}+x_{i+2}=x_ix_{i+1},\quad i\in\Z.
\end{equation}
Three consecutive entries determine every entry in both directions.
The inverse in \eqref{eq:map} follows directly from this recurrence.
Also,
\begin{align*}
 K(T(x,y,z))
 &=y^2+z^2+(yz-x)^2-yz(yz-x)\\
 &=x^2+y^2+z^2-xyz=K(x,y,z),
\end{align*}
which proves the invariance in \eqref{eq:invariant}.

\begin{lemma}[Reversal]\label{lem:reversal}
If $(x_i)_{i\in\Z}$ satisfies \eqref{eq:recurrence}, then so does
$y_i=x_{-i}$. Periodicity and the maximum absolute coordinate are
preserved under this reversal.
\end{lemma}
\begin{proof}
The second equality in \eqref{eq:recurrence}, with index $-i-1$,
is
$x_{-i-2}+x_{-i+1}=x_{-i-1}x_{-i}$.
This is exactly $y_{i-1}+y_{i+2}=y_iy_{i+1}$, with the summands
and factors interchanged. The remaining assertions follow from the
bijective reindexing of $\Z$.
\end{proof}

\begin{proposition}[The displayed orbits]\label{prop:itineraries}
The sets in Table~\ref{tab:families} are periodic orbits with exactly
the stated least periods, heights, and levels. They are pairwise
disjoint, including across all values of $m\geq1$.
\end{proposition}
\begin{proof}
Read each of the following scalar words cyclically:
\begin{align}
 A_m &: (m,0,0,-m,0,0),\label{eq:wordA}\\
 B_m &: (-1,m,-1,1-m),\label{eq:wordB}\\
 C_m &: (1,m,1,m-1,-1,-m,1,1-m,1,-m,-1,m-1).
 \label{eq:wordC}
\end{align}
Each consecutive triple is a point, and each cyclic shift is one
application of $T$. Direct substitution in
\eqref{eq:recurrence} verifies the first two words. For example, in
the four-letter word the first newly generated entry is
$m(-1)-(-1)=1-m$, the next is $(-1)(1-m)-m=-1$, and the
remaining two are $m$ and $-1$. For the twelve-letter word, the
complete list of triples is shown in Table~\ref{tab:twelve}; applying
$T$ to index $j$ gives index $j+1$ modulo twelve.

\begin{table}[htbp]
\centering
\begin{tabular}{@{}rl@{\qquad}rl@{}}
\toprule
$j$ & $T^j(1,m,1)$ & $j$ & $T^j(1,m,1)$\\
\midrule
$0$ & $(1,m,1)$ & $6$ & $(1,1-m,1)$\\
$1$ & $(m,1,m-1)$ & $7$ & $(1-m,1,-m)$\\
$2$ & $(1,m-1,-1)$ & $8$ & $(1,-m,-1)$\\
$3$ & $(m-1,-1,-m)$ & $9$ & $(-m,-1,m-1)$\\
$4$ & $(-1,-m,1)$ & $10$ & $(-1,m-1,1)$\\
$5$ & $(-m,1,1-m)$ & $11$ & $(m-1,1,m)$\\
\bottomrule
\end{tabular}
\caption{The twelve consecutive triples of $C_m$. The numbering refers
to ordinary iterates of $T$, with no sign quotient or change of clock.}
\label{tab:twelve}
\end{table}

The positive value $m$ occurs once in \eqref{eq:wordA}. Therefore
that scalar word cannot be a repetition of a shorter word, and the
least period is six. For $m\geq2$ the same unique occurrence of the
positive value $m$ in \eqref{eq:wordB} and \eqref{eq:wordC}
proves least periods four and twelve. This argument needs a separate
check at $m=1$, where the value $1$ occurs more than once. In that
case the four triples of $B_1$ are
\begin{equation}\label{eq:B1}
 (-1,1,-1),\quad(1,-1,0),\quad(-1,0,-1),\quad(0,-1,1),
\end{equation}
which are distinct. The twelve triples of Table~\ref{tab:twelve}
at $m=1$ are likewise pairwise distinct: those containing a zero
have different ordered positions or signs, and the four without a
zero are
$(1,1,1),(-1,-1,1),(1,-1,-1),(-1,1,-1)$.
Hence the stated least periods hold at this endpoint as well.

For the four points in $E\sqcup D$, each coordinate has modulus two
and their product is positive. Thus $yz=2x$ at each point, so that
$T(x,y,z)=(y,z,x)$. This gives the fixed point $E$ and the
three-cycle $D$. The origin is also fixed.

Since $|m-1|\leq m$ for $m\geq1$, the scalar words have heights
exactly $m$. At the representatives, direct evaluation gives
\begin{equation}\label{eq:family-levels}
 K(m,0,0)=m^2,\qquad
 K(-1,m,-1)=K(1,m,1)=m^2-m+2.
\end{equation}
The special levels follow by substitution. The three infinite families
have distinct least periods, so no orbit belongs to two different
families. Within a family, different $m$ give different heights.
The special orbits have least periods one or three and are disjoint
from those families; $O$ and $E$ are different points. This proves
all the assertions.
\end{proof}

The proposition establishes inclusion of the right-hand side of
\eqref{eq:classification} in the periodic set. It does not yet
exclude other periods. That exclusion is the next section's purpose.
